\ifdefined\iscombined
	\lecturetitle{Interaction-Oriented Architectures}
\else
\documentclass{article}
\usepackage{listings}
\usepackage{amsthm}
\usepackage{comment}
\usepackage{amsmath}
\usepackage{amsfonts}
\usepackage{sectsty}
\usepackage{hyperref}
\usepackage{fancyhdr}
\usepackage[top=1in,bottom=1in,left=1in,right=1in]{geometry}
\usepackage{float}
\usepackage{graphicx}
\usepackage{tabularx}
\usepackage{mathtools}

\date{
	\vspace{-.75em}
	{\large \today}
}

\title{
	\vspace*{-1.5em}
	{\Huge Lecture 15} \\
	\vspace{-1em}
	\rule{\textwidth/4}{.01em} \\
	\vspace{-.25em}
	{\Large Interaction-Oriented Architectures}
	\vspace{-.75em}
}

\author{
	{\large Matthew Farah}
}

\hypersetup{
	colorlinks=true,
	linkcolor=blue,
	filecolor=magenta,
	urlcolor=blue,
	citecolor=blue,
	anchorcolor=blue
}

\fancyhead{}
\fancyhead[L]{\today}
\fancyhead[R]{Lecture 15 - Interaction-Oriented Architectures}
\sectionfont{\fontsize{12}{12}\bfseries}
\subsectionfont{\fontsize{10}{12}\normalfont}
\subsubsectionfont{\fontsize{10}{12}\normalfont}

\begin{document}

\maketitle
\pagestyle{fancy}

\fi

\begin{itemize}
	\item The two major interaction-oriented architectures covered are:
	\begin{enumerate}
		\item Presentation, Abstraction, Control (PAC).
		\item Model, View, Control (MVC).
		\begin{itemize}
			\item MVC actually has two closely related variants we will also cover:
			\begin{enumerate}
				\item MVC-I.
				\item MVC-II.
			\end{enumerate}
		\end{itemize}
	\end{enumerate}
	\item PAC and MVC are closely related.
	\begin{itemize}
		\item The presentation module is analogous to the view module.
		\item The abstraction module is analogous to the model module.
	\end{itemize}
	\item PAC and MVC are differentiated by their overall structure.
	\begin{itemize}
		\item MVC is composed of a single model, view, and controller element.
		\item PAC is composed of a hierarchy of presentation, abstraction, and control elements.
		\begin{itemize}
			\item An `agent' is a group of these three elements in the hierarchy working together.
		\end{itemize}
	\end{itemize}
	\item PAC agents only communicate with one another by their control elements.
	\item A new agent is required for each distinct source of data handled by the system.
	\begin{itemize}
		\item Data sources are distinct if they describe different quantities or are sent to the system at different frequencies.
	\end{itemize}
	\item There are benefits and drawbacks to using a PAC architecture.
	\begin{itemize}
		\item Benefits:
		\begin{enumerate}
			\item Comparatively better at handling multiple, different kinds of data.
			\item Comparatively better at handling data being received at varying frequencies.
		\end{enumerate}
		\item Drawbacks:
		\begin{enumerate}
			\item If a PAC hierarchy is deep, the top-level controller takes comparatively long to communicate with the bottom-level controller.
			\item If the PAC hierarchy is shallow, the top-level controller has high coupling.
		\end{enumerate}
	\end{itemize}
	\item An effective implementation of a PAC architecture involves properly balancing the trade-offs associated with different hierarchical organizations of agents.
	\item MVC is primarily used for simpler, web-based systems.
	\begin{itemize}
		\item MVC emphasizes the system's look-and-feel.
	\end{itemize}
	\item Certain software principles may be occasionally ignored if violating it is more advantageous than adhering.
	\item MVC-II is the usual variant of MVC discussed above.
	\item MVC-I is a simpler variant of MVC composed of two elements:
	\begin{enumerate}
		\item A model module.
		\begin{itemize}
			\item Forms the system's `back-end'.
		\end{itemize}
		\item A single view/control module.
		\begin{itemize}
			\item Forms the system's `front-end'.
		\end{itemize}
	\end{enumerate}
	\item MVC-I violates the separation of concerns principle.
	\begin{itemize}
		\item For simple systems with a minimal need for a separate controller module, this is worth it.
	\end{itemize}
	\item The subscribe/notify pattern is used to implement communication between elements in MVC modules.
	\begin{itemize}
		\item In the subscribe/notify patterns, elements can be either:
		\begin{enumerate}
			\item Subjects.
			\item Observers.
		\end{enumerate}
		\item Observers can subscribe to subjects to be notified of state changes.
		\item Observers can unsubscribe to subjects to stop receiving such notifications.
		\item Subjects track all subscribed observers.
		\item Subjects have a special method to notify all subscribed observers.
	\end{itemize}
	\item In the MVC-I architecture:
	\begin{itemize}
		\item The back-end is the subject.
		\item The front-end is the observer.
	\end{itemize}
	\item In the MVC-II architecture:
	\begin{itemize}
		\item The model element is the subject.
		\item The controller and view elements are both observers of the model element.
	\end{itemize}
\end{itemize}

\ifdefined\iscombined
	\newpage
\else
	\end{document}
\fi
